\chapter{}

A scuola Pietro motteggiava i più vicini di banco con la sua ilarità nervosa; li costringeva a dargli retta, li chiamava con soprannomi faceti, li offendeva se non gli davano retta. E anche quando tutti tacevano, nè meno udiva la voce dell’insegnante; quantunque qualche risposta dei compagni gli arrivasse agli orecchi con un rammarico strano.

Stava per prendere la licenza elementare, ed era il più grande e il meno bravo; e i seminaristi lo canzonavano.

Qualche volta, dopo aver cercato di comprendere, si sforzava di badare a tutta la lezione rimanente; e sentiva quasi gusto ad aumentare la disistima di tutti, benchè se ne compiangesse. Quando era stato attento, usciva con la mente quasi stravolta, con un peso dentro le tempie, incapace di mettersi a studiare; stanco sfinito; senza aver fatto nulla: lasciava un libro e ne prendeva un altro, lasciava anche questo e non leggeva; non s’accorgeva nè meno più d’averli dinanzi.

Allora, si divertiva al movimento e al vocìo della trattoria.

Del resto avrebbe dovuto imparare le sue lezioni e scrivere dinanzi agli avventori meno ricchi che desinavano a una tavola lunga, sopra alla quale ciascuno di loro distendeva un piccolo tovagliolo: lungo i solchi delle piegature, si raccoglievano le briciole del pane, e Pietro le mangiava a pizzichi.

Questi avventori, divenuti amici di Domenico e di Anna, lo facevano ridere con le loro burlette dicendogli:

— Che vuoi affaticarti gli occhi? Vai a ruzzare.

Ma Anna si alzava dalla sua poltrona posta nell’angolo più oscuro della stanza, dietro un paravento di legno con un’apertura rotonda da cui poteva sorvegliare i camerieri e la cantiniera, per dire:

— Lo lascino stare!

Poi, rideva anche lei.

D’estate, quando tirava un poco di vento, si vedeva uscire dalla finestra aperta tutto il fumo delle pipe e dei sigari; e allora gli avventori si toglievano la giubba; mentre, d’inverno, si passavano uno scaldino.

Burlavano tra loro, portandosi via il pane o le frutta. Quando qualcuno bestemmiava troppo, Anna impallidiva e lo guardava in faccia. Egli rimaneva con la parola in bocca e tutti gli altri tacevano; e la conversazione era cambiata.

— La bestemmia non sta bene. Avete tempo fuori di qui! Per la strada!

Quegli arrossiva:

— Ieri il rimprovero non toccò a me! Non è vero, padrona?

Era una risata spontanea. Ed Anna pensava subito ad un’altra cosa.

Allora qualcuno proponeva:

— Venga a darci da bere. Ma non di quello annacquato. Non ci castighi!

Chi aveva ancora un poco di vino, vuotava il bicchiere riposandolo con gli altri nel mezzo della tavola. Anna, fattosi portare un fiasco, domandava:

— Quanto ne vuole lei?

— Un soldo.

— Io due soldi....

Adamo metteva il bicchiere contro l’aria:

— È piovuto in cantina anche oggi!

Quando passava un avventore delle altre sale, si chetavano alla meglio e lo seguivano con lo sguardo.

— È il tale.

Qualche volta, cantavano. Ma Domenico usciva dalla cucina tenendo un ramaiolo di brodo. Tutti alzavano le mani:

— Fermo! Fermo! Ce ne andiamo!

Gli alterchi erano radi; e, quando avvenivano, l’amicizia era rotta per poco tempo. Di solito, non s’insultavano direttamente; ma uno alla volta, a vicenda, si rivolgevano agli altri esponendo la cosa come un racconto; da prima a bassa voce, poi con veemenza e con bestemmie, battendo i pugni, alzandosi da sedere.

Quasi, le mani dei contendenti si toccavano; allora, qualcuno diceva:

— È vergogna! Anche per chi ci sente!

Anna non si teneva più; e la sfilata delle bestemmie era interrotta, finalmente, da un grosso boccone inghiottito.

Adamo, con piccole nervosità da femmina avvezzata male, quando diceva a Domenico che lo servisse bene, quasi si raccomandava. Dopo averlo guardato in viso, si volgeva da una parte, aspettando, sempre con la paura che parlassero male di lui in cucina; poi, assaggiata due o tre volte la pietanza, se era a modo suo respirava meglio, sputacchiava e si decideva a mangiare. E, tornatagli la gaiezza, era primo lui a svegliare Giacomino, mettendogli una buccia di mela nel collo. Anziano, basso e corpulento, con i baffi sempre in bocca, cambiava d’umore come un ragazzo. Anzi, chiedeva scusa dell’inurbanità del momento prima, battendo insieme le dita sopra il tovagliolo, tamburellandole, con la testa in avanti e bassa. Si stropicciava le guance con il dorso della mano, silenzioso, con il sigaro in bocca, biascicandolo e facendolo girare tra le labbra. Era capace di mettersi ad ascoltare una lunga conversazione fatta nella stanza accanto; per dirne, con una frase sola o con un sospiro, la sua opinione. E, se per caso gli avessero risposto, si rifaceva pensoso, fumando a boccate più lunghe.

Giacomino, anche mangiando, appoggiava la testa alla mano, tirandosi con le dita i capelli vicini alla nuca.

Bibe metteva il mento sopra il pugno chiuso, in proda alla tavola, e stava così con gli occhi giù, divertendosi ad ascoltare, senza veder nessuno; e allora alzava, una per volta e piano, le punte dei piedi, battendole in tempo; finchè qualcuno, presolo per i capelli ricciuti, non gli facesse volger la testa.

— Dio! Mi fate male! Che divertimento c’è?

— Hai sonno, bestia?

— Poco no.

E raccontava perchè non aveva avuto tempo di dormire abbastanza. E sorrideva, tra il sonno.

Volevano sempre gli stessi posti: Adamo in un angolo, perchè così sputava a piacere; Giacomino sotto la finestra; Bibe il più giovine, sul canapè: perchè ci si tirava in dietro a modo suo, magari addormentandocisi quando non gli davano fastidio.

Si riabbottonavano i calzoni, si riagganciavano gli scheggiali, sputavano, s’urtavano, si scapaccionavano, si tiravano i baffi e pagavano il conto andando, a uno per volta, dinanzi al bugigattolo di Anna.

E Pino? Pino, il vecchio barrocciaio di Poggibonsi, era il più povero. Gridava, per ridere:

— C’è posto anche per me?

Tutti glielo facevano, non per cortesia, ma perchè lo credevano pieno di pulci. Egli se ne avvedeva, ma non osava dir niente: brontolava un poco tra sè; e, siccome dovunque era trattato così, non se la prendeva.

— Mezzo posto mi basta a me. Non sono un signore io! Ah, come mi dolgono le ossa!

Un occhio non gli voleva stare aperto, e le palpebre battevano insieme come fanno quelle delle civette. Girava quell’altro occhio per tutta la stanza, lentamente; ricominciando sempre da capo. Si guardava bene le mani, per far capire agli altri che aveva pensato a lavarsele; e in fatti se l’era lavate nel secchio del suo cavallo mezzo stronco come le stanghine del barroccio, rinforzate con parecchie avvolte di funicella e di filo di ferro. Quanto tempo gli faceva perdere quel lavoro riaccomodato tutti i giorni!

Si stropicciava gli occhi con un dito, con il viso ridente senza sapere perchè: la sua bocca, con quel sorriso, pareva larga il doppio.

— Ridete voi, eh, boia! Che avete rubato oggi? Si piglia la roba delle commissioni e poi dice che l’ha persa per la strada.

— Io? Oh, poverino! Una volta lo facevo così, ma ora no.

Strascicava la voce con un accento, che sembrava sincero benchè malizioso. E poi:

— Ho due figliole, a casa, da maritare! Son belle da vero, a dirvela in un orecchio. Ma la mia moglie è già ridotta come una balla di cenci unti, che non si piglierebbero nè meno in mano. Ci ho quelle due figliole, povere bambine! O che devo fare io per loro?

Tutta la sua fisonomia pigliava una bontà umile ma ostinata; e, cosa strana, le sue guancie, tra il pelo della barba rada, erano delicate come quelle di una donna.

Egli non ordinava, ma Domenico gli sceglieva tra la roba del giorno innanzi e gliene faceva un piatto solo. Lo pigliava per la tesa del cappello: quasi gli ci faceva battere il naso:

— Senti come ti ho servito?

— Sì, avete ragione, è stantia, ma non puzza tanto.

Adamo e Giacomino gli buttavano fette di pane o mezze frutta. Egli, senza guardarli in faccia, se le radunava più vicino, quasi avesse voluto metterle sotto il tondo del piatto, con ambedue le mani.

— Oh, oggi sto meglio!

Salutava con molto rispetto Anna, aspettando che gli rispondesse: e, certo, non si sarebbe messo a sedere prima. Tanto che Anna, quando se n’era dimenticata, doveva dirgli:

— Mettetevi a sedere!

— Ah, mi ci posso mettere? Credevo dì dar noia oggi! Sono tanto stanco!

E aspettava, tenendo le mani insieme.

Da Pietro si faceva rispiegare, quasi una volta al mese, che cosa erano le due oleografie delle pareti. Pietro saliva in piedi su la panca, per non staccarle. Ma Pino diceva:

— Me le metta più vicino! Se sapesse, Pietrino, come mi bruciano gli occhi! Qualche volta ho paura d’accecare.

Una era la Battaglia di Adua e un’altra I fattori dell’unità italiana. E tenendolo, dopo, per una manica:

— Non dia retta al babbo: studi. Me ne intendo!

Pietro, allora, senza sapere perchè, lo accarezzava.

D’inverno, quando era tutto infreddolito e bagnato, con il bavero della giubba fino alla cima degli orecchi, con il cappello su gli occhi, Pietro gli si faceva subito incontro; e, senza parlargli, gli metteva il viso tanto vicino, che Domenico lo tirava in dietro per il collo.

Morì presto; e nessuno se ne accorse.
